\documentclass[11pt,a4paper]{article}
% Common preamble for the hint sheets (\input from each file)
\usepackage[utf8]{inputenc}
\usepackage[T1]{fontenc}
\usepackage{lmodern}
\usepackage{amsmath,amssymb,amsthm,mathtools}
\usepackage[a4paper,margin=2.3cm]{geometry}
\usepackage{enumitem}
\usepackage{xcolor}
\usepackage[hidelinks]{hyperref}
\renewcommand{\leq}{\leqslant}
\renewcommand{\le}{\leqslant}
\renewcommand{\geq}{\geqslant}
\renewcommand{\ge}{\geqslant}
\setlength{\parindent}{0pt}
\setlength{\parskip}{0.35em}
% \hintstitle{Sheet title}{Class work / Homework}
\newcommand{\hintstitle}[2]{\begin{center}{\LARGE\bfseries Hints: #1\par}\vspace{0.5em}{\large #2\par}\end{center}
\noindent\textit{The hints to each problem are cumulative. Read one, try again, and open the next only if you are still stuck.}\par\medskip\hrule\medskip}
\newcommand{\problem}[1]{\par\bigskip{\large\bfseries Problem #1}\par\nopagebreak\smallskip}
\newcommand{\subpart}[1]{\par\smallskip\textbf{(#1)}\par\nopagebreak}
\newcommand{\hint}[1]{\par\smallskip\noindent\textcolor{black!70}{\textbf{Hint #1.}}\enspace}
\newcommand{\note}[1]{\par\smallskip\noindent\textit{Note.} #1\par}
\providecommand{\lt}{<}
\providecommand{\gt}{>}
\providecommand{\tightlist}{\setlength{\itemsep}{0pt}\setlength{\parskip}{0pt}}
\emergencystretch=3em

\begin{document}
\hintstitle{Methods of proof -- 4. Invariants, monovariants and double counting}{Class work}

\problem{1}
\hint{1} Treat the empty cell as a sixteenth tile, so a position is a permutation $\pi$ of the 16 cells. Every move is a transposition of the empty cell with a neighbouring tile. What does a move do to $\operatorname{sgn}\pi$, and what does it do to the colour of the empty cell in the chessboard colouring?
\hint{2} Each move changes both: the sign flips and the empty cell changes colour. So $I=\operatorname{sgn}\pi\cdot(-1)^{c}$, where $c$ is the colour (equivalently, the parity of the taxicab distance of the empty cell from its home), is invariant.
\hint{3} Swapping two tiles with the empty cell at home changes $\operatorname{sgn}\pi$ but not $c$, so it changes $I$: unreachable. General principle: the reachable configurations from $x$ form the orbit of $x$ under the group (or semigroup) generated by the moves; any function constant on these orbits, typically $x\mapsto\chi(g)$ for a homomorphism $\chi$ (here to $\{\pm1\}$) that is trivial on all moves, separates orbits.

\problem{2}
\hint{1} The map $(a,b)\mapsto\bigl(\frac{a+b}{\sqrt2},\frac{a-b}{\sqrt2}\bigr)$ is given by an orthogonal $2\times2$ matrix. Which simple symmetric functional of the numbers on the board does it preserve?
\hint{2} The sum of squares is invariant, so a multiset with a different sum of squares is unreachable. It is not sufficient. With two numbers write the board as a point $(a,b)$ of the plane; the move is the reflection in the line at angle $22.5^\circ$, and since the order of the two numbers is free, the swap (reflection in the line at $45^\circ$) is available too. Their product is the rotation by $45^\circ$, so they generate the symmetry group of a regular octagon (order $16$), and the reachable multisets are exactly $\{\pm a,\pm b\}$ and $\bigl\{\pm\frac{a+b}{\sqrt2},\pm\frac{a-b}{\sqrt2}\bigr\}$ with all choices of signs: at most $8$ multisets on the circle $x^2+y^2=a^2+b^2$, although the invariant allows the whole circle.
\hint{3} For $(a,b)\mapsto(a+b,ab)$: the new pair is the root pair of $(x-s)(x-p)$, where $s=a+b$, $p=ab$, so on $(s,p)$ the operation acts as $(s,p)\mapsto(s+p,sp)$; no symmetric polynomial is preserved. Use a monovariant instead: the sum grows by $ab$, so on positive numbers the sum strictly increases, hence no configuration occurs twice, and only multisets with larger sum (and positive entries) can be reached. (Positivity matters: a $0$ on the board is frozen, since $(a,0)\mapsto(a,0)$, and a negative product decreases the sum.)

\problem{3}
\hint{1} For positive integers $a\ge b>0$ the sum $a+b$ (or the larger number) strictly decreases at each step and is a positive integer, so the process stops, namely when one number becomes $0$.
\hint{2} The invariant is $\gcd(a,b)=\gcd(a-b,b)$; at the end it is the remaining non-zero number. Grouping consecutive subtractions of the same $b$ gives division with remainder: the Euclidean algorithm.
\hint{3} Complexity: subtraction alone can take $a$ steps (e.g.\ $(a,1)$). With division, if $r_{k}=q r_{k+1}+r_{k+2}$, then $r_k\ge r_{k+1}+r_{k+2}$, so the remainders dominate Fibonacci numbers backwards; the number of division steps is $O(\log b)$ (Lam\'e: at most about $\log_\varphi b+1$; also $r_{k+2}<r_k/2$).
\hint{4} One continuous analogue: for the gradient flow $x'=-\nabla V(x)$ the energy $V$ is a Lyapunov function, $V(x(t))=V(x(0))-\int_0^t|x'(s)|^2\,ds$. If $V$ is bounded below, $\int_0^\infty|x'|^2<\infty$, which is the continuous form of ``a bounded monovariant cannot decrease forever''. (Note also that Euclid on reals $a,b$ terminates iff $a/b\in\mathbb Q$.)

\problem{4}
\hint{1} $\sum_v\deg v=2|E|$: count pairs $(v,e)$ with $v$ an endpoint of $e$. For the 4-cycle bound count the ``cherries'' $(v,\{x,y\})$ with $x\ne y$ both adjacent to $v$.
\hint{2} By vertices $v$ there are $\sum_v\binom{d_v}{2}$ cherries; by pairs $\{x,y\}$ there are at most $\binom n2$, because without a $C_4$ two vertices have at most one common neighbour. Hence $\sum_v\binom{d_v}2\le\binom n2$.
\hint{3} By convexity (Jensen), $n\binom{d}{2}\le\binom n2$ with $d=2|E|/n$, i.e.\ $d^2-d-(n-1)\le0$, so $d\le\frac{1+\sqrt{4n-3}}2$ and $|E|=\frac{nd}2\le\frac12\bigl(1+\sqrt{4n-3}\bigr)\cdot\frac n2$.
\hint{4} Oddtown: sets $A_1,\dots,A_m\subset\{1,\dots,n\}$ with odd sizes and even pairwise intersections. Their characteristic vectors over $\mathbb F_2$ have Gram matrix $I_m$, which has rank $m$, so the vectors are independent and $m\le n$. Comparison: double counting bounds a quantity by averaging over incidences (and the $C_4$ bound is asymptotically sharp); for Oddtown a counting argument gives only weaker bounds, while the rank gives $m\le n$ at once, which is sharp (take singletons).

\problem{5}
\hint{1} Encode $S$ by a generating function: $|S|$ is the coefficient of $x^{3860}$ in $f(x)^{2018}$, where $f(x)=x+x^2+x^3+x^4+x^5+x^6+x^{10}$.
\hint{2} All coefficients are non-negative, so for every $x>0$ we have $|S|\,x^{3860}\le f(x)^{2018}$, i.e.\ $|S|\le f(x)^{2018}/x^{3860}$. Choose $x$ cleverly.
\hint{3} Take $x=\frac12$: $f(\frac12)=\frac{63}{64}+\frac1{1024}=\frac{1009}{1024}=\frac{2018}{2048}$, which gives exactly $|S|\le2^{3860}\bigl(\frac{2018}{2048}\bigr)^{2018}$.

\end{document}
